\documentclass[11pt]{article}
\usepackage[T1]{fontenc}
\usepackage{textcomp}
\usepackage[margin=0.75in]{geometry}
\usepackage{amsmath,amssymb,amsthm}
\usepackage{array,booktabs}
\usepackage{hyperref}
\setlength{\parindent}{0pt}
\newtheorem{theorem}{Theorem}
\theoremstyle{definition}
\newtheorem{definition}{Definition}
\title{Flow Proof Artifact: prop-47-pythagoras}
\author{Generated by \texttt{flow doc proof}}
\date{Flow Proof Book}
\begin{document}
\maketitle
In a right-angled triangle, the square on the hypotenuse equals the sum of squares on the legs.\\[0.5em]
\textbf{Source.} euclid — Elements, Book I, Proposition 47\\[0.5em]
\bigskip
\noindent{\large \textbf{Derived fact 1.} \textit{In a right-angled triangle, the square on the hypotenuse equals the sum of squares on the legs} \hfill the Euclidean plane \cdot Euclid Book I \cdot Proposition 47: square on hypotenuse equals sum of squares on legs}\par
\smallskip\noindent\textit{Coordinate: the Euclidean plane \cdot Euclid Book I \cdot Proposition 47: square on hypotenuse equals sum of squares on legs}\par
\smallskip\noindent\textit{Source: euclid — Elements, Book I, Proposition 47}\par
\smallskip\noindent\textit{Built on: proposition 4: side-angle-side congruence, for Book I of the Elements in on the Euclidean plane, proposition 31: interior angles on one side sum to two right angles, for Book I of the Elements in on the Euclidean plane, proposition 46: describe a square on a given straight line, for Book I of the Elements in on the Euclidean plane}\par
\medskip
\noindent\textbf{Goal.} In a right-angled triangle, the square on the hypotenuse equals the sum of squares on the legs\par
\noindent$\displaystyle c^{2} = a^{2} + b^{2}$\par
\medskip
\noindent
\begin{tabular}{@{} >{\raggedright\arraybackslash}p{0.06\textwidth} >{\raggedright\arraybackslash}p{0.47\textwidth} >{\raggedleft\arraybackslash}p{0.41\textwidth} @{}}
\textbf{\#} & \textbf{Proof} & \textbf{Mathematics} \\
\midrule
\textbf{1.} & We prove that proposition 47: square on hypotenuse equals sum of squares on legs for Book I of the Elements in the Euclidean plane. & \\[0.45em]
\textbf{2.} & Let square\_on\_AB = square\_constructed\_on\_AB\_by\_Proposition\_46. & \\[0.45em]
\textbf{3.} & Let square\_on\_AC = square\_constructed\_on\_AC. & \\[0.45em]
\textbf{4.} & Let square\_on\_BC = square\_on\_hypotenuse. & \\[0.45em]
\textbf{5.} & Let line\_through\_C\_parallel\_to\_BF = parallel\_through\_C. & \\[0.45em]
\textbf{6.} & We invoke the derived fact: triangle\_ABC\_right\_angled\_at\_A. & \\[0.45em]
\textbf{7.} & We invoke the axiom governing Book I of the Elements in the Euclidean plane: proposition 4: side-angle-side congruence, for Book I of the Elements in on the Euclidean plane. & \\[0.45em]
\textbf{8.} & We invoke the derived fact governing Book I of the Elements in the Euclidean plane: proposition 31: interior angles on one side sum to two right angles, for Book I of the Elements in on the Euclidean plane. & \\[0.45em]
\textbf{9.} & We invoke the derived fact governing Book I of the Elements in the Euclidean plane: proposition 46: describe a square on a given straight line, for Book I of the Elements in on the Euclidean plane. & \\[0.45em]
\textbf{10.} & From step 2, step 3, step 4, step 5, step 6, step 7, step 8, and step 9, we can deduce that area square on BC equals area square on AB plus area square on AC. & $\displaystyle \text{area square on BC} = \text{area square on AB} + \text{area square on AC}$ \\[0.45em]
\textbf{11.} & We can deduce that c times c equals a times a plus b times b. Hence proven. & $\displaystyle c^{2} = a^{2} + b^{2}$ \\[0.45em]
\bottomrule
\end{tabular}
\label{thm:Geometry-Euclid Book I-proposition_47_square_on_hypotenuse_equals_sum_of_squares_on_legs}
\par\medskip\hrule
\end{document}
